\chapter*{LỜI MỞ ĐẦU}
\addcontentsline{toc}{chapter}{LỜI MỞ ĐẦU}

\section*{1. Lý do chọn đề tài}

Trong bối cảnh toàn cầu hóa và sự bùng nổ của ngành thương mại điện tử, ngành Logistics và Quản trị Chuỗi cung ứng (\textit{Supply Chain Management}) đã trở thành hạ tầng xương sống quyết định năng lực cạnh tranh của nền kinh tế quốc gia. Việc triển khai các dự án logistics quy mô lớn --- chẳng hạn như xây dựng trung tâm phân phối đa vùng, hiện đại hóa hệ thống kho bãi tự động hay phát triển hệ thống phần mềm quản lý vận tải --- đòi hỏi sự phối hợp nhịp nhàng giữa hàng trăm công việc phụ thuộc phức tạp và hàng loạt các loại tài nguyên hạn chế \cite{christopher2016}.

Tuy nhiên, việc quản trị tiến độ và chi phí trong các dự án logistics thực tế đang đối mặt với những thách thức ngặt nghèo:
\begin{enumerate}
    \item \textbf{Cấu trúc phụ thuộc mạng lưới đa chiều và hiệu ứng Lan truyền trễ (Delay Propagation):} Sự chậm trễ của một công việc ở thượng nguồn (ví dụ: thông quan hải quan hoặc kiểm thử phần mềm) có thể nhanh chóng khuếch đại thành sự đứt gãy tiến độ toàn bộ hạ nguồn, làm bùng nổ chi phí lưu kho và phạt vi phạm hợp đồng \cite{pmbok2017}.
    \item \textbf{Bài toán Lập lịch Ràng buộc Tài nguyên (RCPSP) và Đánh đổi Thời gian -- Chi phí (Time-Cost Trade-off):} Trong thực tế, tài nguyên nhân lực chuyên môn và thiết bị luôn bị giới hạn ngặt nghèo. Khi xảy ra xung đột tài nguyên, các nhà quản lý buộc phải thực hiện điều phối tài nguyên (\textit{Resource Leveling}) hoặc rút ngắn tiến độ (\textit{Crashing}) bằng cách làm thêm giờ hoặc thuê nhà thầu phụ, làm thay đổi phức tạp cấu trúc chi phí dự án \cite{kolisch1997psplib}.
    \item \textbf{Hạn chế của các phần mềm quản lý truyền thống:} Các giải pháp thương mại hiện nay chủ yếu dựa trên mô hình bảng phẳng hoặc các thuật toán tất định CPM truyền thống, gặp bẫy hiệu năng $N+1$ khi truy vấn mạng lưới đồ thị quy mô lớn, đồng thời thiếu khả năng học tập tự động các rủi ro bất định từ dữ liệu lịch sử.
\end{enumerate}

Sự xuất hiện của công nghệ \textbf{Trí tuệ Nhân tạo Đồ thị (Graph AI)} --- đặc biệt là kiến trúc \textbf{Heterogeneous Graph Transformer (HGT)} \cite{hu2020hgt} kết hợp \textbf{Học tự giám sát (Self-Supervised Learning -- SSL)} \cite{he2022mae} --- mở ra cơ hội đột phá giúp mô hình hóa tự nhiên cấu trúc đa chiều của dự án logistics. Đồng thời, việc kết hợp Graph AI với các bộ tối ưu hóa số nguyên (Constraint Programming / MILP) và thuật toán tiến hóa (Genetic Algorithms) tạo tiền đề xây dựng một giải pháp quản trị thông minh, đáp ứng thời gian thực. Vì những lý do cấp thiết đó, đồ án lựa chọn nghiên cứu đề tài: \textbf{"Hệ thống Quản trị Tiến độ và Chi phí Dự án Logistics dựa trên AI Đồ thị Dị thể và Tối ưu hóa Lai đa mục tiêu (GLPO)"}.

\section*{2. Mục tiêu nghiên cứu}

Đồ án đặt ra mục tiêu tổng quát là nghiên cứu, thiết kế và phát triển thành công hệ thống \textbf{GLPO (Logistics Project Progress and Cost Management System)} nhằm tự động hóa quy trình quản lý tiến độ, dự báo rủi ro lan truyền trễ và tối ưu hóa cân bằng đa mục tiêu (Thời gian -- Chi phí -- Tài nguyên) cho các dự án logistics. Để đạt được mục tiêu tổng quát, 5 mục tiêu cụ thể được triển khai:

\begin{enumerate}
    \item \textbf{Xây dựng Động cơ Truy vấn Đồ thị Đệ quy (Recursive CTE Engine):} Thiết kế cấu trúc lưu trữ và giải pháp duyệt sơ đồ mạng PDM trực tiếp trong cơ sở dữ liệu PostgreSQL, triệt tiêu bẫy hiệu năng $N+1$ Query và đảm bảo tốc độ phản hồi thời gian thực ($< 100$ ms).
    \item \textbf{Phát triển Phân hệ Trí tuệ Nhân tạo Đồ thị (Heterogeneous Graph AI):} Biểu diễn dự án dưới dạng Đồ thị Dị thể ($V_{\text{Task}}$, $V_{\text{Resource}}$, $V_{\text{Shift}}$), huấn luyện mô hình HGT tích hợp cơ chế SSL Masked Autoencoder (MAE) để dự báo độ trễ công việc và trích xuất Trọng số Chú ý (Attention Weights) giải thích các nút găng rủi ro.
    \item \textbf{Xây dựng Bộ Tối ưu hóa Lai Đa mục tiêu (Hybrid MILP-GA Engine):} Phối hợp giữa giải thuật Quy hoạch tuyến tính số nguyên hỗn hợp (MILP / CP-SAT) và Thuật toán Di truyền (GA) với cơ chế "Tiêm hạt giống ưu tú" (\textit{Elite Seeding}), giải quyết trọn vẹn bài toán RCPSP và Crashing tiến độ.
    \item \textbf{Hiện thực hóa Mô hình Tài chính Phân ly 38 loại Chi phí (G1--G7):} Xây dựng bộ công cụ phân tích cấu trúc ngân sách đa tầng, hỗ trợ quản lý chỉ số EVM ($PV, EV, AC, CPI, SPI$) và thực hiện chiến lược Crashing có mục tiêu.
    \item \textbf{Đóng gói và Triển khai Đóng chuẩn (Benchmark Implementation):} Triển khai toàn bộ hệ thống theo kiến trúc vi dịch vụ Microservices (React, FastAPI, PyTorch Geometric, Docker), tiến hành kiểm thử và đối chuẩn toàn diện trên bộ dữ liệu chuẩn quốc tế \textbf{UniSA MFET 3008 / MFET 5040}.
\end{enumerate}

\section*{3. Đối tượng và phạm vi nghiên cứu}

\textbf{Đối tượng nghiên cứu} của đồ án bao gồm:
\begin{itemize}
    \item Cấu trúc mạng lưới công việc AON DAG, quan hệ phụ thuộc tiến trình PDM (FS, SS, FF) và các chỉ số thời gian CPM ($ES, EF, LS, LF, Slack$).
    \item Định ngạch tiêu thụ tài nguyên nhân lực chuyên môn hạn chế ($R_1 \rightarrow R_8$) và phương thái tiêu thụ tài nguyên theo thời gian.
    \item Mô hình toán học của bài toán lập lịch ràng buộc tài nguyên (RCPSP), bài toán rút ngắn tiến độ (Crashing) và cơ cấu 38 loại chi phí logistics thuộc 7 nhóm chỉ số ($G_1 \rightarrow G_7$).
    \item Kiến trúc mạng nơ-ron đồ thị dị thể HGT, cơ chế học tự giám sát SSL MAE và bộ giải tối ưu lai MILP-GA.
\end{itemize}

\textbf{Phạm vi dữ liệu và đối chuẩn:} Nghiên cứu tập trung đánh giá định lượng trên bộ dữ liệu Case Study chuẩn quốc tế \textbf{MFET 3008 / MFET 5040} (\textit{Parcel Transport, Storage and Database System -- University of South Australia}) bao gồm $21$ nút công việc cốt lõi, $26$ đường liên kết phụ thuộc, $8$ chuyên môn nhân sự hạn chế và ngân sách chi phí trực tiếp cơ sở $49.510$ USD.

\textbf{Phạm vi triển khai kỹ thuật:} Đồ án xây dựng ứng dụng web hoàn chỉnh theo kiến trúc Microservices đóng gói Container bằng Docker Compose, không bao gồm việc tích hợp trực tiếp vào phần mềm phần cứng IoT của xe tải trên thực địa.

\section*{4. Phương pháp nghiên cứu}

Đồ án ứng dụng kết hợp phương pháp nghiên cứu lý thuyết, mô hình hóa toán học và thực nghiệm kỹ thuật:
\begin{enumerate}
    \item \textbf{Phương pháp Mô hình hóa Toán học và Tối ưu hóa:} Sử dụng Lý thuyết Sơ đồ mạng AON để xây dựng đồ thị DAG, thiết lập mô hình bài toán quy hoạch toán học cho RCPSP và Crashing. Áp dụng kỹ thuật nới lỏng MILP kết hợp với thuật toán tiến hóa di truyền GA để tìm tập nghiệm Pareto.
    \item \textbf{Phương pháp Học sâu Đồ thị (Graph Deep Learning):} Xây dựng mạng HGT trên thư viện PyTorch Geometric, áp dụng kỹ thuật tự giám sát Masked Autoencoder (MAE) với hàm suy hao Smooth L1 Loss để học đại diện đồ thị mà không cần nhãn thủ công.
    \item \textbf{Phương pháp Kỹ thuật Phần mềm (Software Engineering):} Áp dụng quy tắc thiết kế vi dịch vụ (Microservices), xây dựng RESTful API với FastAPI, tối ưu hóa truy vấn CSDL với PostgreSQL Recursive CTE và thiết kế giao diện tương tác với React-Flow.
    \item \textbf{Phương pháp Đánh giá Thực nghiệm (Empirical Benchmark):} Thực hiện đo đạc độ trễ truy vấn (\textit{Query Latency}), dung lượng bộ nhớ tiêu thụ (\textit{RAM/VRAM Footprint}), tốc độ hội tụ di truyền, sai số hồi quy AI (MAE, RMSE, MAPE, $R^2$) và kiểm định giả thuyết thống kê (\textit{Wilcoxon test}).
\end{enumerate}

\section*{5. Cấu trúc của đồ án}

Nội dung chính của đồ án được tổ chức thành 4 chương logic:

\textbf{Chương 1: Tổng quan lý thuyết và Cơ sở nghiên cứu} --- Trình bày đặc thù dự án logistics (mô hình SCOR, hiệu ứng Bullwhip), lý thuyết sơ đồ mạng AON, CPM, bài toán RCPSP, bài toán Crashing, phương pháp EVM, cơ sở lý thuyết về Mạng nơ-ron Đồ thị Dị thể (HGT) và Học tự giám sát (SSL MAE).

\textbf{Chương 2: Kiến trúc Hệ thống và Phương pháp luận đề xuất} --- Mô tả chi tiết kiến trúc tổng thể Microservices, lược đồ CSDL PostgreSQL, thuật toán truy vấn Recursive CTE, đường ống tiền xử lý ETL Tensor, thuật toán Tối ưu hóa Lai MILP-GA, pipeline huấn luyện AI HGT-SSL và hạ tầng đóng gói Docker.

\textbf{Chương 3: Kết quả Thực nghiệm và Thảo luận} --- Trình bày toàn bộ số liệu đối chuẩn thực nghiệm định lượng trên Case Study UniSA MFET 3008/5040, bao gồm hiệu năng truy vấn CTE, mức tiêu thụ bộ nhớ RAM/VRAM, hiệu quả làm mịn tài nguyên (\textit{Resource Leveling}), kết quả Crashing tiến độ, độ chính xác mô hình AI (MAE, RMSE, MAPE, $R^2$), khả năng giải thích bằng Attention Heatmap và bàn luận chuyên sâu.

\textbf{Chương 4: Kết luận và Hướng phát triển} --- Tổng kết các đóng góp khoa học và thực tiễn của đề tài, phân tích hạn chế nội tại và đề xuất 3 hướng phát triển trọng tâm (Thuật toán Di truyền Thích nghi AGA, kiến trúc Graphormer Advanced Attention và mở rộng chuỗi dữ liệu thời gian thực SaaS/ERP).